% GENERATED FILE. Edit canonical lessons, then run npm run generate:books.
% canonical-source-hash: fnv1a64:b7e48fab656ad95e
% canonical-lessons: ZH-C413-dengji-board, ZH-C413-fei, ZH-W413-fei-fly, ZH-C413-jizhang, ZH-C413-jiangluo, ZH-C413-feikuai, ZH-R413-review-408-413

\chapter{Up in the Air}
\label{ch:zh-up-in-the-air}
\hlchaptermodality{413}

\begin{chapteropening}
\textbf{By the end of this chapter:} \emph{I can board a plane, say how high it flies, talk about the captain and landing, and say time flies.}

Complete the last lesson of chapter 413: I can board a plane, say how high it flies, talk about the captain and landing, and say time flies.
\end{chapteropening}

\section[dēngjī]{dēngjī — to board a plane}

\label{lesson:ZH-C413-dengji-board}



\begin{quote}
Before the new one: say the Mandarin for to come in to shore, then the Mandarin for a passenger ship.
\end{quote}

\subsection*{You'll want to know: dēngjī}
\emph{dēngjī} — \textquotedblleft{}to board a plane\textquotedblright{}.

\begin{quote}
\emph{dēngjī yǐqián wǒmen zuò zài zhèr děng} — Until we board, we'll sit here and wait.
\end{quote}

Its characters come later; for now it is a word for the ear.

\begin{cousinweb}
\emph{dēngjī} is built from \emph{dēng} 'to go up' and \zh{机} 'a plane', both inherited from Middle Chinese, \emph{tong} and \emph{kj+j} (Guangyun, in Baxter's transcription), which glosses the first 'to rise' and, after an older dictionary, 'to get up into a carriage'.
\end{cousinweb}

\begin{grammarlens}[title={What you've built}]
One more word for getting around.
\end{grammarlens}

\subsection*{Your turn}
\begin{itemize}
\raggedright
  \item \emph{Say it:} \emph{dēngjī}
  \item \emph{Say it:} \emph{dēngjī}, once more
  \item \emph{Say it:} \emph{dēngjī yǐqián wǒmen zuò zài zhèr děng} — Until we board, we'll sit here and wait.
  \item \emph{Recall:} say the Mandarin for to come in to shore, then the Mandarin for a passenger ship, then say \emph{dēngjī} again
\end{itemize}

\begin{tcolorbox}[breakable,colback=teal!4,colframe=teal!35!black,title={Before you move on}]
What does dēngjī mean? (\textquotedblleft{}To board a plane\textquotedblright{}.) Say it once more.
\end{tcolorbox}
\section[\texorpdfstring{\zh{飞} (fēi)}{飞 (fēi)}]{\zh{飞} (fēi) — to fly}

\label{lesson:ZH-C413-fei}



\begin{quote}
Before the new one: say the Mandarin for a passenger ship, then the Mandarin for to board a plane.
\end{quote}

\subsection*{You'll want to know: \zh{飞}}
\textbf{\zh{飞}} — \emph{fēi} — \textquotedblleft{}to fly\textquotedblright{}.

\begin{quote}
\textbf{\zh{飞机飞得很高}} \emph{fēijī fēi de hěn gāo} — The plane is flying very high.
\end{quote}

The new character, in pieces: \zh{飞} — three strokes: a bar that turns and slants down to the right, ending in a hook up; then a short stroke falling to the left and a dot, near the top on the right. The next lesson puts it in your hand.

\begin{cousinweb}
\textbf{\zh{飞}} (\emph{fēi}) is inherited from Middle Chinese \emph{pj+j} (Guangyun, in Baxter's transcription), glossed 'to fly, to soar'.
\end{cousinweb}

\begin{grammarlens}[title={What you've built}]
One more word for getting around.
\end{grammarlens}

\subsection*{Your turn}
\begin{itemize}
\raggedright
  \item \emph{Say it:} \emph{fēi}
  \item \emph{Say it:} \emph{fēi}, once more
  \item \emph{Say it:} \emph{fēijī fēi de hěn gāo} — The plane is flying very high.
  \item \emph{Recall:} say the Mandarin for a passenger ship, then the Mandarin for to board a plane, then say \emph{fēi} again
\end{itemize}

\begin{tcolorbox}[breakable,colback=teal!4,colframe=teal!35!black,title={Before you move on}]
What does \zh{飞} mean? (\textquotedblleft{}To fly\textquotedblright{}.) Say it once more.
\end{tcolorbox}
\section[\texorpdfstring{\zh{飞} (fēi)}{飞 (fēi)}]{\zh{飞} — a hooked turn, then two small strokes}

\label{lesson:ZH-W413-fei-fly}



\begin{quote}
The word first, then the character inside it.

\begin{itemize}
\raggedright
  \item \emph{Recall:} say \emph{fēi} — \textbf{R1}, one lesson back, before it decays
\end{itemize}

\begin{quote}
\zh{机} \zh{起}
\end{quote}

\emph{Read it:} each one aloud, with nothing said first — \emph{jī}, \emph{qǐ}
\end{quote}

\subsection*{Script — a hooked turn, then two small strokes}
\hlblockfigure{\detokenize{figures/ZH-W413-fei-fly-filmstrip.pdf}}{How \zh{飞} is written, stroke by stroke}

\begin{quote}
\zh{飞}
\end{quote}

\textbf{One piece, three strokes.}

\textbf{First, a bar to the right that turns and slants down to the right}, then flicks up at the end: one stroke, without lifting the pen.

\textbf{Then a short stroke falling to the left}, on the right of the turn.

\textbf{Last, a dot} under it, falling to the right.

One long stroke and two short ones: three.

\subsection*{Your turn}
\begin{itemize}
\raggedright
  \item \emph{Write it:} \zh{飞}, all three, with the model in front of you
  \item \emph{Cover:} write it again from memory, and count the pieces, not the strokes
\end{itemize}

\begin{tcolorbox}[breakable,colback=teal!4,colframe=teal!35!black,title={Before you move on}]
How many strokes in \zh{飞}? (\textbf{Three}.) Which stroke is the longest? (\textbf{The first, the bar that turns and hooks}.)

Source: \href{https://raw.githubusercontent.com/chanind/hanzi-writer-data/68d10a4b21150cae5e1ebbd223eed289cf32d90c/data/\%E9\%A3\%9E.json}{Hanzi Writer Data \zh{飞}, pinned stroke paths and medians}.
\end{tcolorbox}
\section[\texorpdfstring{\zh{机长} (jīzhǎng)}{机长 (jīzhǎng)}]{\zh{机长} (jīzhǎng) — plane's captain; pilot in command}

\label{lesson:ZH-C413-jizhang}



\begin{quote}
Before the new one: say the Mandarin for to board a plane, then the Mandarin for to fly.
\end{quote}

\subsection*{You'll want to know: \zh{机长}}
\textbf{\zh{机长}} — \emph{jīzhǎng} — \textquotedblleft{}plane's captain; pilot in command\textquotedblright{}.

\begin{quote}
\textbf{\zh{机长说飞机马上要到了}} \emph{jīzhǎng shuō fēijī mǎshàng yào dào le} — The captain says we'll be there soon.
\end{quote}

\begin{cousinweb}
\textbf{\zh{机长}} (\emph{jīzhǎng}) is built from \zh{机} 'a plane' and \zh{长} 'the head', both inherited from Middle Chinese, \emph{kj+j} and \emph{trjangX} (Guangyun, in Baxter's transcription).
\end{cousinweb}

\begin{grammarlens}[title={What you've built}]
One more word for getting around.
\end{grammarlens}

\subsection*{Your turn}
\begin{itemize}
\raggedright
  \item \emph{Say it:} \emph{jīzhǎng}
  \item \emph{Say it:} \emph{jīzhǎng}, once more
  \item \emph{Say it:} \emph{jīzhǎng shuō fēijī mǎshàng yào dào le} — The captain says we'll be there soon.
  \item \emph{Recall:} say the Mandarin for to board a plane, then the Mandarin for to fly, then say \emph{jīzhǎng} again
\end{itemize}

\begin{tcolorbox}[breakable,colback=teal!4,colframe=teal!35!black,title={Before you move on}]
What does \zh{机长} mean? (\textquotedblleft{}Plane's captain; pilot in command\textquotedblright{}.) Say it once more.
\end{tcolorbox}
\section[jiàngluò]{jiàngluò — to land; to come down onto the runway}

\label{lesson:ZH-C413-jiangluo}



\begin{quote}
Before the new one: say the Mandarin for to fly, then the Mandarin for the pilot in command.
\end{quote}

\subsection*{You'll want to know: jiàngluò}
\emph{jiàngluò} — \textquotedblleft{}to land; to come down onto the runway\textquotedblright{}.

\begin{quote}
\emph{fēijī jiàngluò yǐhòu wǒmen zuò dìtiě qù bīnguǎn} — Once the plane has landed, we'll take the metro to the hotel.
\end{quote}

Its characters come later; for now it is a word for the ear.

\begin{cousinweb}
\emph{jiàngluò} is built from \emph{jiàng} 'to come down' and \emph{luò} 'to fall', both inherited from Middle Chinese, \emph{kaewngH} and \emph{lak} (Guangyun, in Baxter's transcription), which glosses the first 'down, to come down' and the second as leaves falling from a tree.
\end{cousinweb}

\begin{grammarlens}[title={What you've built}]
One more word for getting around.
\end{grammarlens}

\subsection*{Your turn}
\begin{itemize}
\raggedright
  \item \emph{Say it:} \emph{jiàngluò}
  \item \emph{Say it:} \emph{jiàngluò}, once more
  \item \emph{Say it:} \emph{fēijī jiàngluò yǐhòu wǒmen zuò dìtiě qù bīnguǎn} — Once the plane has landed, we'll take the metro to the hotel.
  \item \emph{Recall:} say the Mandarin for to fly, then the Mandarin for the pilot in command, then say \emph{jiàngluò} again
\end{itemize}

\begin{tcolorbox}[breakable,colback=teal!4,colframe=teal!35!black,title={Before you move on}]
What does jiàngluò mean? (\textquotedblleft{}To land; to come down onto the runway\textquotedblright{}.) Say it once more.
\end{tcolorbox}
\section[\texorpdfstring{\zh{飞快} (fēikuài)}{飞快 (fēikuài)}]{\zh{飞快} (fēikuài) — very fast; as fast as flying}

\label{lesson:ZH-C413-feikuai}



\begin{quote}
Before the new one: say the Mandarin for the pilot in command, then the Mandarin for to come down onto the runway.
\end{quote}

\subsection*{You'll want to know: \zh{飞快}}
\textbf{\zh{飞快}} — \emph{fēikuài} — \textquotedblleft{}very fast; as fast as flying\textquotedblright{}.

\begin{quote}
\textbf{\zh{时间过得飞快}} \emph{shíjiān guò de fēikuài} — Time flies.
\end{quote}

\begin{cousinweb}
\textbf{\zh{飞快}} (\emph{fēikuài}) is built from \zh{飞} 'to fly' and \zh{快} 'fast', both inherited from Middle Chinese, \emph{pj+j} and \emph{khwaejH} (Guangyun, in Baxter's transcription), which glosses the second 'to one's liking'; the sense 'fast' is not given there.
\end{cousinweb}

\begin{grammarlens}[title={What you've built}]
One more word for getting around.
\end{grammarlens}

\subsection*{Your turn}
\begin{itemize}
\raggedright
  \item \emph{Say it:} \emph{fēikuài}
  \item \emph{Say it:} \emph{fēikuài}, once more
  \item \emph{Say it:} \emph{shíjiān guò de fēikuài} — Time flies.
  \item \emph{Recall:} say the Mandarin for the pilot in command, then the Mandarin for to come down onto the runway, then say \emph{fēikuài} again
\end{itemize}

\begin{tcolorbox}[breakable,colback=teal!4,colframe=teal!35!black,title={Before you move on}]
What does \zh{飞快} mean? (\textquotedblleft{}Very fast; as fast as flying\textquotedblright{}.) Say it once more.
\end{tcolorbox}
\section[(dialogue)]{Looking back: from By Bus to Up in the Air}

\label{lesson:ZH-R413-review-408-413}



\begin{quote}
Say the Mandarin for a bus stop sign and public transport.
\end{quote}

\subsection*{Your turn}
\begin{itemize}
\raggedright
  \item \emph{Say it:} a bus stop sign, public transport, to drop coins in the box, a long-distance coach, a starting point
  \item \emph{Say it:} the rush hour, the traffic on the roads, to open to traffic, a ticket seller, to take a taxi
  \item \emph{Say it:} long-distance, the fare for a ride, to pass through, to launch a new transport line, to hit the road
  \item \emph{Say it:} a pier for boats, a ship's captain, to set sail, a ferry across the water, a life jacket
  \item \emph{Say it:} a harbour for ships, the ship's deck, a ship's cabin, to come in to shore, a passenger ship
  \item \emph{Say it:} to board a plane, to fly, the pilot in command, to come down onto the runway, as fast as flying
  \item \emph{Read it:} \zh{交} \zh{通} \zh{船} \zh{飞} — each one aloud, then a word that holds it
\end{itemize}

\begin{tcolorbox}[breakable,colback=teal!4,colframe=teal!35!black,title={Before you move on}]
A bus stop sign? (\emph{zhànpái}.) Public transport? (\textbf{\zh{公交}}, \emph{gōngjiāo}.) To drop coins in the box? (\emph{tóubì}.) A long-distance coach? (\textbf{\zh{客车}}, \emph{kèchē}.) A starting point? (\textbf{\zh{起点}}, \emph{qǐdiǎn}.) The rush hour? (\emph{gāofēng}.) The traffic on the roads? (\textbf{\zh{交通}}, \emph{jiāotōng}.) To open to traffic? (\textbf{\zh{通车}}, \emph{tōngchē}.) A ticket seller? (\emph{shòupiàoyuán}.) To take a taxi? (\textbf{\zh{打车}}, \emph{dǎchē}.) Long-distance? (\emph{chángtú}.) The fare for a ride? (\emph{chēfèi}.) To pass through? (\textbf{\zh{通过}}, \emph{tōngguò}.) To launch a new transport line? (\textbf{\zh{开通}}, \emph{kāitōng}.) To hit the road? (\textbf{\zh{上路}}, \emph{shànglù}.) A pier for boats? (\emph{mǎtou}.) A ship's captain? (\textbf{\zh{船长}}, \emph{chuánzhǎng}.) To set sail? (\textbf{\zh{开船}}, \emph{kāichuán}.) A ferry across the water? (\emph{dùchuán}.) A life jacket? (\textbf{\zh{救生衣}}, \emph{jiùshēngyī}.) A harbour for ships? (\emph{gǎngkǒu}.) The ship's deck? (\emph{jiǎbǎn}.) A ship's cabin? (\emph{chuáncāng}.) To come in to shore? (\emph{kào'àn}.) A passenger ship? (\textbf{\zh{客船}}, \emph{kèchuán}.) To board a plane? (\emph{dēngjī}.) To fly? (\textbf{\zh{飞}}, \emph{fēi}.) The pilot in command? (\textbf{\zh{机长}}, \emph{jīzhǎng}.) To come down onto the runway? (\emph{jiàngluò}.) As fast as flying? (\textbf{\zh{飞快}}, \emph{fēikuài}.)
\end{tcolorbox}
